\documentclass[11pt,a4paper]{article}
\usepackage[T1]{fontenc}
\usepackage{isabelle,isabellesym}

\usepackage{amsmath}
\usepackage{amssymb}
\usepackage{stmaryrd}

% this should be the last package used
\usepackage{pdfsetup}

\urlstyle{rm}
\isabellestyle{it}

\begin{document}

\title{The Space of Continuous Paths, and Tightness from Increment Moments}
\author{Dmitriy Traytel}
\maketitle

\begin{abstract}
  Weak convergence of laws of continuous processes needs a Polish state space
  and a tightness criterion.  This entry supplies both.

  The space $C([0,T],\,'b)$ of continuous paths with the supremum metric is
  built as a metric space in the set-based form of the
  \texttt{Standard\_Borel\_Spaces} entry, and shown separable and complete.
  H\"older balls in it are compact, which is Arzel\`a--Ascoli in the form
  needed later.  The path map of a stochastic process is measurable into it,
  so a process has a law on path space; these laws are consistent across
  horizons.  The continuous mapping theorem and two portmanteau facts are
  proved for weak convergence of such laws: a closed set of full measure
  survives a weak limit, and an open set of positive mass survives a weak
  perturbation.

  The tightness criterion is a moment estimate.  For a martingale whose
  compensator grows at a controlled rate, second and fourth moments of the
  increments are bounded on a partition; a dyadic chaining argument turns the
  fourth-moment bound into a tail bound for the modulus of continuity, and
  H\"older interpolation converts that into a bound on a H\"older seminorm.
  Since H\"older balls are compact, the family of laws is tight, hence
  relatively compact by Prokhorov's theorem from the
  \texttt{Levy\_Prokhorov\_Metric} entry.

  The fourth-moment estimate is proved without the Burkholder--Davis--Gundy
  inequality, by iterating Cauchy--Schwarz and the tower property; the
  resulting constant is smaller than the one the usual route
  gives~\cite{LaiShkolnikovSoner}.  Billingsley~\cite{Billingsley} is the
  textbook reference for the tightness argument.
\end{abstract}

\tableofcontents

\parindent 0pt\parskip 0.5ex

\input{session}

\bibliographystyle{abbrv}
\bibliography{root}

\end{document}
